\hypertarget{perimeter-and-consolidation-state}{%
\section{Perimeter and consolidation state}\label{method-groups-fx-methods-perimeter-and-consolidation-state}}

The acquisition, control and currency methods below use the inspected IFRS 3, IFRS 10 and IAS 21 provisions in the retained 2023 EU text \citep{eu20231803}. The acquisition with goodwill is a new example; it does not change the preceding zero-goodwill case.

ControlAssessment records investor/investee, relevant activities, rights and their substantive/protective character, variable returns, ability to use power, principal/agent analysis, effective interval and evidence. IFRS 10.5--9 requires power, returns and their connection, with reassessment on changed facts. An 80\% share count is evidence to analyze, not an automatic substitute for control. Unknown control yields separate conditional perimeters and unresolved group results; no entity disappears merely to make reconciliation easy.

PerimeterVersion binds control intervals, ownership classes, report cutoffs, functional/presentation currency, policy-alignment adjustments and the exact source entity books. ConsolidationAdjustment is a distinct group book with matched counterparty/event references; legal-entity journals remain intact. The procedure is: select time-compatible entity books; apply supported policy alignment; translate them; add like balances/movements; eliminate investment against acquisition equity; eliminate intragroup balances, income/cost and unrealised gains; allocate profit/OCI to owner/NCI; reconcile all statements and movement notes. IFRS 10.19--23/B86--B96 support these steps. Reporting-date differences need B92--B93's conditions and intervening adjustments; a generic nearest-date join is invalid.

\hypertarget{acquisition-and-ownership-case}{%
\section{Acquisition and ownership case}\label{method-groups-fx-methods-acquisition-and-ownership-case}}

Assume a business combination, no common-control transaction, parent obtaining control and 80\% of ordinary residual interests for cash 96. Acquiree identifiable assets fair value 160 and liabilities 60 give net 100. There is no prior interest, contingent consideration, tax difference, separately unrecognized intangible, transaction cost or impairment; those absences are synthetic facts. Select proportionate-share NCI under IFRS 3.19: NCI 20 and goodwill 96+20\ensuremath{-}100=16 under .32. In parent books debit investment 96/credit cash 96. In the consolidation overlay, debit acquiree acquisition equity 100 and goodwill 16, credit investment 96 and NCI 20. Only post-acquisition income enters group performance. Asset measurement exceptions, deferred tax and a bargain gain reassessment under .36 would change the calculation; goodwill cannot absorb an unexplained missing balance.

If post-acquisition net profit is 20, all accrued in retained earnings with no dividend/OCI, NCI increases 4 to 24 and parent's attributable profit is 16. Parent buys another 10\% for 15 while retaining control: transfer 12 from NCI (10\% of current net 120), credit cash 15 and debit parent equity 3. Remaining NCI 12; goodwill 16 unchanged; no new profit. IFRS 10.23/B96 makes this an equity transaction. This example has equal ordinary rights; preference rights need their actual allocation and cannot use a simple ownership fraction.

In a separate alternative starting before that additional purchase, sell all 80\% for 130 when net assets 120, goodwill 16, NCI 24 and OCI zero. On loss of control, group gain=130+24\ensuremath{-}(120+16)=18. Derecognize individual assets/liabilities and goodwill, remove NCI, recognize consideration and gain. The separate parent book's investment gain is not added again. If 20\% is retained and 60\% sold for 100 with retained fair value 35, group gain=100+35+24\ensuremath{-}136=23; recognize the retained investment 35. Its subsequent associate/financial-asset method needs separately established significant-influence/classification facts. IFRS 10.B98--B99 also requires the appropriate prior-OCI treatment; this case assumes none.

\hypertarget{intercompany-inventory-and-tax}{%
\section{Intercompany inventory and tax}\label{method-groups-fx-methods-intercompany-inventory-and-tax}}

Assume parent's external inventory cost 60 is sold to its controlled subsidiary for 100 on credit, all still held at report date, tax ignored in the base case. Parent records AR 100/revenue 100 and COGS 60/inventory 60; subsidiary inventory 100/ AP 100. Consolidation debits AP 100/credits AR 100 and debits revenue 100/credits COGS 60 and inventory 40. Group inventory 60, revenue and COGS zero, and no group profit arises. When subsidiary later sells externally for 120, group COGS is 60 and profit 60, not 20; release the previously eliminated 40 in the continuing adjustment roll-forward. Reversing a closing overlay and recreating a new one must not recognize the profit twice.

If the purchaser's independently supported tax base is 100, group carrying amount 60 creates deductible difference 40. At assumed applicable 25\% and with all IAS 12 DTA conditions satisfied, recognize group DTA 10/deferred-tax benefit 10; reverse on outside sale. The original tax expense and cash obligations in the seller's tax jurisdiction remain. Different rates, recoverability or tax-law transfer rules change the answer; tax is not simply eliminated with group profit. For an upstream subsidiary sale, attribute elimination and related tax effects to the appropriate subsidiary/NCI performance; downstream parent's profit elimination does not mechanically reduce NCI by the same fraction.

\hypertarget{functional-currency-versus-group-presentation}{%
\section{Functional currency versus group presentation}\label{method-groups-fx-methods-functional-currency-versus-group-presentation}}

EntityCurrencyState names the functional currency and its evidenced rationale, foreign-currency nominal balances, monetary/nonmonetary classification, historical or fair-value measurement dates, rates and sources. IAS 21.21--23 and .28 require transaction-date translation and closing-rate remeasurement of monetary items; historical-cost nonmonetary assets retain transaction-date rates. A functional currency change is prospective under .35--37. Hyperinflation, lack of exchangeability and hedge accounting require separate applicable methods; do not import the historical .26 text into a later amended period without checking it.

Example: buy goods for foreign 100 at functional 1.2 per foreign unit; debit inventory 120/credit payable 120. Closing rate 1.3 gives payable 130 and FX expense 10, while inventory remains 120 before any NRV test. Pay at 1.25: debit payable 130/ credit cash 125/credit FX gain 5. Cumulative expense 5 matches 125 cash less 120 asset cost. Translating the inventory again to 130 would be wrong for this declared basis.

Group presentation is a different operation under IAS 21.39--41. A subsidiary has opening net assets 100, equity 100, no debt or distributions, annual profit 20 entirely earned at rate 1.25, and closing assets 120. Opening historical rate 1.2 and closing 1.3 give translated assets 156, contributed/opening equity 120 and profit 25; translation OCI 11 reconciles the difference. It is explained as 100\ensuremath{\times}(1.3\ensuremath{-}1.2)+20\ensuremath{\times}(1.3\ensuremath{-}1.25)=11, not a free balancing plug. If ownership is 80\%, parent/NCI shares of profit are 20/5 and OCI 8.8/2.2 under the stated equal rights. Use transaction rates; an average is acceptable only as a justified approximation, not when significant fluctuation invalidates it. Goodwill and fair-value adjustments follow IAS 21.47's foreign-operation treatment, rather than freezing them automatically in the parent's currency.

Group reports retain the translation reserve by entity, ownership and event. Loss of control requires the appropriate reserve reclassification under IAS 21.48 and IFRS 10.B99; a no-control-loss change transfers the relevant attribution in equity under applicable conditions. Cash translation effects reconcile cash separately from operating/investing/financing movements. An intercompany monetary balance's exchange result cannot simply be eliminated as if no currency exposure existed; IAS 21.45 and net-investment .32 require the substantive classification.
